\documentclass[10pt]{article}
\usepackage[a4paper,margin=1in]{geometry}
\usepackage{amsmath}
\usepackage{booktabs}
\usepackage{hyperref}
\usepackage{enumitem}

\title{Volatility Inference with SDEs and Bayesian Filtering:\\
A Short Updated Report}
\author{}
\date{April 2026}

\begin{document}
\maketitle

\begin{abstract}
This report summarizes an updated volatility inference study for cryptocurrency data using stochastic differential equation (SDE) state models with Bayesian filtering. Two pipelines were tested: (i) lag-aware rolling volatility estimation and (ii) instantaneous volatility estimation using scaled absolute log-return as a proxy. After tightening the evaluation protocol to reduce optimistic bias, rolling-volatility performance became more realistic. In this setting, Kalman/Particle filtering remained competitive, while benchmark GARCH models performed better on some assets. The instantaneous proxy setup remained unstable under extreme crypto fluctuations.
\end{abstract}

\section{Problem Statement}
Crypto volatility is hard to model because of rapid sentiment shifts, liquidity shocks, leverage cascades, and frequent regime changes. The goal is online volatility estimation with fair out-of-sample comparison against GARCH baselines.

\section{Methodology}
\subsection{Pipeline 2 (Primary): Rolling Volatility, Lag-Aware}
Latent volatility is modeled as a mean-reverting state:
\begin{equation}
\sigma_{t+1}=\sigma_t+\kappa(\theta-\sigma_t)\Delta t+\xi\eta_t,\quad \eta_t\sim\mathcal{N}(0,\Delta t).
\end{equation}
The scored target is rolling volatility $y_t^{(roll)}=\mathrm{std}(r_{t-w+1:t})$, while assimilation input is lagged $z_t^{(DA)}=y_{t-L}^{(roll)}$.

Kalman prediction/update (scalar):
\begin{align}
\hat x_t^- &= \hat x_{t-1}+\kappa(\theta-\hat x_{t-1})\Delta t,\quad P_t^- = P_{t-1}+Q,\\
K_t &= \frac{P_t^-}{P_t^-+R},\quad \hat x_t=\hat x_t^-+K_t(y_t-\hat x_t^-),\quad P_t=(1-K_t)P_t^-.
\end{align}
Particle filtering uses bootstrap propagation and Gaussian likelihood weighting.

\subsection{Pipeline 1 (Secondary): Instantaneous Proxy}
Instantaneous proxy target:
\begin{equation}
y_t^{(inst)}=\left|\log\frac{S_t}{S_{t-1}}\right|\sqrt{\frac{\pi}{2}}.
\end{equation}
This was treated as an exploratory proxy for latent volatility under the same filtering framework.

\section{Evaluation Protocol}
To reduce overestimation:
\begin{itemize}[leftmargin=1.2em]
  \item score DA \textbf{prior} signals (not posterior updates),
  \item exclude burn-in samples,
  \item in rolling mode, enforce explicit lag between DA input and scored target,
  \item score GARCH on post-warmup forecasts only.
\end{itemize}

\section{Results (BTC/ETH Snapshot)}
\subsection{Rolling (Current, Lag-Aware: target \texttt{rolling\_vol\_20\_lag20})}
\begin{table}[h]
\centering
\begin{tabular}{lccc}
\toprule
Symbol & Kalman $R^2$ & Particle $R^2$ & GARCH(2,2) $R^2$ \\
\midrule
BTC & 0.1416 & 0.1044 & 0.3295 \\
ETH & 0.2375 & 0.2387 & 0.1708 \\
\bottomrule
\end{tabular}
\end{table}

\subsection{Rolling (Legacy, no explicit lag)}
Legacy scores were much higher (Kalman $R^2\approx0.992$--$0.993$ for BTC/ETH), but this reflected smooth target dynamics and online tracking alignment rather than robust forecasting realism.

\subsection{Instantaneous Proxy (Current)}
For BTC/ETH, Kalman and Particle Heston scores remained negative or near zero, while short rolling baselines were modestly positive. This indicates that the proxy-based instantaneous setup is currently unreliable in high-fluctuation crypto regimes.

\section{Conclusion}
The main practical outcome is that lag-aware rolling volatility inference is feasible and competitive, but not uniformly superior to GARCH across assets. The previous near-perfect rolling scores were optimistic under a less strict setup. Instantaneous volatility estimation from log-return proxies did not generalize well in this study.

\section{Future Work}
Current work uses mid-price-derived returns for simplicity. Next steps are:
\begin{itemize}[leftmargin=1.2em]
  \item sentiment/news features,
  \item order-book and flow-based signals,
  \item multidimensional state-space models for regime-aware volatility dynamics.
\end{itemize}

\end{document}
